% Source: ME12_v1.html
\documentclass[11pt]{article}
\usepackage[margin=1in]{geometry}
\usepackage{amsmath}
\usepackage{amssymb}
\usepackage{siunitx}
\usepackage{enumitem}

\title{ME12: Fluid Mechanics}
\author{}
\date{}

\begin{document}
\maketitle

\section*{ME12: Fluid Mechanics}

\begin{enumerate}[leftmargin=*,label=\textbf{Question \arabic*.}]
\item \hfill \textit{/ 10 pts}

Calculate the density of a spherical ball that has mass 454 kg and radius of 1 m.



𝜌 = \rule{2cm}{0.4pt} kg/m$^{3}$

\textit{Correct answer: } 108.3846

\item \hfill \textit{/ 10 pts}

Estimate the density of an alloy made from 4.3 kg of silver and 2.3 kg of gold. The density of silver is 10,490 kg/m$^{3}$ and the density of gold is 19,300 kg/m$^{3}$.



𝜌 = \rule{2cm}{0.4pt} kg/m$^{3}$

\textit{Correct answer: } 12,474.3617

\item \hfill \textit{/ 10 pts}

A cube of mass 40 kg and side 19 cm is placed on a smooth, flat surface in a vacuum on a planet where the acceleration of gravity is 2.9 m/s$^{2}$. Calculate the pressure on the surface beneath the cube.



\emph{p} =\rule{2cm}{0.4pt}Pa

\textit{Correct answer: } 3,213.2964

\item \hfill \textit{/ 10 pts}

Consider the barometer in the figure and calculate the ambient pressure in kPa if \emph{h}$_{1}$ = 3.9 cm, \emph{h}$_{2}$ = 31 cm, and the fluid density is 6,282 kg/m$^{3}$.

[Image: barometer01.png]



$p_0=\_\_\_\_\_\;\mathrm{kPa}$

\textit{Correct answer: } 16.6837

\item \hfill \textit{/ 10 pts}

On a far planet the local acceleration due to gravity is 8.3 m/s$^{2}$ and the atmospheric pressure is 80,512 Pa. Calculate the pressure 7.4 m below the surface of a fluid having density 936 kg/m$^{3}$.



\emph{p} = \rule{2cm}{0.4pt} Pa

\textit{Correct answer: } 138,001.12

\item \hfill \textit{/ 10 pts}

The figure shows a glass cylinder of radius 5.1 cm, partially filled with oil of density 0.81 × 10$^{3}$ kg/m$^{3}$. Floating freely on top with no oil leaking around the sides, a 2.9-kg wooden disk supports a 0.5-kg potato. Calculate the gauge pressure in psi at 13.2 cm below the surface of the oil. (1 psi = 6.895 kPa)

[Image: potato\_tube.png]



\emph{p} = \rule{2cm}{0.4pt} psi (Note the units!)

\textit{Correct answer: } 0.7434

\item \hfill \textit{/ 10 pts}

Suppose a hydraulic jack has an input piston of cross-sectional area 0.05 m$^{2}$ and an output piston of area 0.17 m$^{2}$. If the output force is 1,243 N, calculate the input force in Newtons.



\emph{F} = \rule{2cm}{0.4pt} N

\textit{Correct answer: } 365.5882

\item \hfill \textit{/ 10 pts}

A hydraulic brake has master cylinder with area 1.7 cm$^{2}$. The brake cylinder has a cross-sectional area of 7.4 cm$^{2}$. The coefficient of friction between the shoe and wheel drum is 0.5. If the wheel has a radius of 45 cm, determine the magnitude of the frictional torque exerted on the wheel around the axle when the driver applies a force of 49 N on the brake pedal. (Note: generally there is a lever system, as well, to amplify the force of the driver's foot!).

[Image: brake2.png]



𝜏 = \rule{2cm}{0.4pt} Nᐧm

\textit{Correct answer: } 47.9912

\item \hfill \textit{/ 10 pts}

An object of mass 7.8 kg and volume 0.3 m$^{3}$ is completely submerged in a tank of fluid with density 609 kg/m$^{3}$. If the acceleration due to gravity is 6.3 m/s$^{2}$at this location (another planet!), calculate the buoyant force \emph{B} on the object.



\emph{B} = \rule{2cm}{0.4pt} N

\textit{Correct answer: } 1,151.01

\item \hfill \textit{/ 10 pts}

A rock is tied to a light cord and suspended in water. Calculate the tension in the cord if the rock's mass and density are 6.4 kg and 3.1 × 10$^{3}$ kg/m$^{3}$.



\emph{T } = \rule{2cm}{0.4pt} N

\textit{Correct answer: } 42.4877

\item \hfill \textit{/ 10 pts}

A pipe of radius 1.1 m narrows to 0.4 m. If the fluid is traveling at 2.9 m/s in the larger radius pipe, calculate the fluid speed in the narrower pipe.



\emph{v }= \rule{2cm}{0.4pt} m/s

\textit{Correct answer: } 21.9312

\item \hfill \textit{/ 10 pts}

A hose with radius 3 cm fills a rectangular pool that is 10 m in width, 3 m in width, and 4 m in depth in 3.5 hours. At what speed does water come from the hose?



\emph{v }= \rule{2cm}{0.4pt} m/s

\textit{Correct answer: } 3.3684

\item \hfill \textit{/ 10 pts}

A large water tank is filled to a depth of 16.3 m. If a leak develops in the side of the tank 8 m above the bottom, at what speed does water shoot out of the leak?



\emph{v }= \rule{2cm}{0.4pt} m/s

\textit{Correct answer: } 12.7546

\item \hfill \textit{/ 10 pts}

A horizontal pipe of radius 0.25 m carries water at 2.3 m/s and pressure of 239 kPa. If the pipe narrows to a radius of 0.11 m, calculate the pressure \emph{p} in kPa in the narrow pipe. The density of water is 1000 kg/m$^{3}$.



\emph{p} = \rule{2cm}{0.4pt} kPa

\textit{Correct answer: } 171.0758

\item \hfill \textit{/ 10 pts}

Which of the following statements is true? Choose the best answer.

\begin{itemize}
  \item The buoyant force is independent of an object’s submerged volume.
  \item For equal masses, the denser mass will have the greater buoyant force in a fluid.
  \item The buoyant force is caused by the weight of the submerged object.
  \item The buoyant force is always greater than the weight of the displaced fluid.
  \item \textbf{[correct]} The buoyant force is a result of increasing pressure with depth.
\end{itemize}

\item \hfill \textit{/ 10 pts}

Two solid metal balls all have the same radius. One is made of aluminum, the other of gold, and their densities are known. If the balls are resting at the bottom of a pool of water, which ball has the largest buoyant force?

\begin{itemize}
  \item \textbf{[correct]} The buoyant forces on the balls are the same.
  \item The aluminum ball has the larger buoyant force.
  \item More information is required.
  \item The gold ball has the larger buoyant force.
\end{itemize}

\end{enumerate}

\end{document}
